% ============================================================
%  Chương VIII – Kết luận
%  ~200 từ · 0.5 trang
% ============================================================
\chapter{Kết luận}
\label{chap:conclusion}

% Ý chính: PKI là "hạ tầng vô hình" — không ai nhìn thấy nhưng mọi người phụ thuộc.
% Không chỉ là vấn đề kỹ thuật mà là vấn đề trust và governance ở quy mô toàn cầu.

% ----------------------------------------------------------
\section{Tổng kết}
\label{sec:8.1}
% Ý chính:
%   - PKI giải quyết bài toán trust ở quy mô internet:
%     ràng buộc (identity ↔ public key) thông qua chữ ký CA
%   - Hành trình logic: symmetric (O(n²) problem)
%     → asymmetric (giải key exchange, tạo authentication problem)
%     → PKI (hoàn chỉnh cả hai)
%   - Web PKI + X.509 phụ trách hàng tỷ kết nối HTTPS mỗi ngày
%   - Thực tế phức tạp hơn lý thuyết: CA có thể bị compromise,
%     revocation không hoàn hảo, governance là thách thức lớn

% ----------------------------------------------------------
\section{Ý nghĩa và triển vọng}
\label{sec:8.2}
% Ý chính:
%   - PKI Việt Nam đang trưởng thành: NEAC Root CA, ký số hành chính,
%     hóa đơn điện tử đang scale nhanh
%   - Ba xu hướng lớn cần theo dõi:
%     post-quantum migration (NIST PQC 2024),
%     automation và short-lived cert (ACME),
%     decentralized identity (W3C DID)
%   - Bài học xuyên suốt: \emph{trust is a sociotechnical problem},
%     không thể giải hoàn toàn bằng kỹ thuật mật mã
